% Part: first-order-logic
% Chapter: completeness
% Section: completeness-theorem

\documentclass[../../../include/open-logic-section]{subfiles}

\begin{document}

\iftag{FOL}
      {\olfileid{fol}{com}{cth}}
      {\olfileid{pl}{com}{cth}}

\olsection{د بشپړتيا قضيه}

\begin{explain}
راځئ خپلې پايلې سره يوځاے کړو: د بشپړتيا قضيې ته رسېږو۔
\end{explain}

\begin{thm}[د بشپړتيا قضيه]
\ollabel{thm:completeness}
فرض کړئ~$\Gamma$ د~!!{sentence}s يو سټ دے۔ که~$\Gamma$ سازګار وي،
نو د صدق وړ دے۔
\end{thm}

\begin{proof}
فرض کړئ~$\Gamma$ سازګار دے۔\iftag{FOL}{ د
  \olref[hen]{lem:henkin} له مخې، يو مشبوع سازګار سټ
  $\Gamma' \supseteq \Gamma$ شته۔}{} د~\olref[lin]{lem:lindenbaum} له
مخې، داسې~$\Gamma^* \supseteq \iftag{FOL}{\Gamma'}{\Gamma}$ شته چې
سازګار او !!{complete} دے۔
\iftag{FOL}{
  ځکه~$\Gamma' \subseteq \Gamma^*$، نو د هرې
  !!{formula}~$!A(x)$ دپاره~$\Gamma^*$ په لاندې بڼه يوه !!a{sentence}
  لري:
  \iftag{prvEx}
        {$\lexists[x][!A(x)] \lif !A(c)$}
        {$\lnot\lforall[x][!A(x)] \lif \lnot !A(c)$}
  او په دې توګه~$\Gamma^*$ مشبوع دے۔ که~$\Gamma$،~$\eq$ نۀ لري،
  نو}{نو} د~\olref[mod]{lem:truth} له مخې،
$\iftag{FOL}{\Sat{M(\Gamma^*)}}{\pSat{v(\Gamma^*)}}{!A}$ هله او
يوازې هله چې~$!A \in \Gamma^*$ وي۔ له دې څخه په تېره بيا دا ترلاسه
کېږي چې د ټولو~$!A \in \Gamma$ دپاره
$\iftag{FOL}{\Sat{M(\Gamma^*)}}{\pSat{v(\Gamma^*)}}{!A}$؛ نو
$\Gamma$ د صدق وړ دے۔\iftag{FOL}{
  که~$\Gamma$،~$\eq$ ولري، نو د~\olref[ide]{lem:truth} له مخې د ټولو
  !!{sentence}s~$!A$ دپاره
  $\Sat{\equivclass{M}{\approx}}{!A}$ هله او يوازې هله چې
  $!A \in \Gamma^*$ وي۔ په تېره بيا، د ټولو~$!A \in \Gamma$ دپاره
  $\Sat{\equivclass{M}{\approx}}{!A}$؛ نو~$\Gamma$ د صدق وړ دے۔}{}
\end{proof}

\begin{cor}[د بشپړتيا قضيه، دويمه بڼه]
\ollabel{cor:completeness}
د ټولو~$\Gamma$ او !!{sentence}s~$!A$ دپاره: که
$\Gamma \Entails !A$ وي، نو~$\Gamma \Proves !A$۔
\end{cor}

\begin{proof}
پام وکړئ چې په~\olref{cor:completeness} او
\olref{thm:completeness} کښې~$\Gamma$ ګانې په کلي ډول کميت ټاکل شوې
دي۔ د دې دپاره چې ځان ګډوډ نۀ کړو،~\olref{thm:completeness} د يوه
بېل متغير په کارولو بيا بيانوو: د~!!{sentence}s د هر سټ~$\Delta$
دپاره، که~$\Delta$ سازګار وي، نو د صدق وړ دے۔ د نقيض عکس له مخې،
که~$\Delta$ د صدق وړ نۀ وي، نو~$\Delta$ ناسازګار دے۔ دا به د لازمې
نتيجې په ثبوت کښې وکاروو۔

فرض کړئ~$\Gamma \Entails !A$۔ بيا د
\olref[syn][sem]{prop:entails-unsat} له مخې
$\Gamma \cup \{\lnot !A\}$ د صدق وړ نۀ دے۔ که
$\Gamma \cup \{\lnot !A\}$ خپل~$\Delta$ ونيسو، د
\olref{thm:completeness} مخکښې بڼه راکوي چې
$\Gamma \cup \{\lnot !A\}$ ناسازګار دے۔ د
\iftag{FOL}{%
  \tagrefs{prfAX/{fol:axd:prv:prop:prov-incons},
    prfSC/{fol:seq:prv:prop:prov-incons},
    prfND/{fol:ntd:prv:prop:prov-incons},
    prfTab/{fol:tab:prv:prop:prov-incons}}}{%
  \tagrefs{prfAX/{pl:axd:prv:prop:prov-incons},
    prfSC/{pl:seq:prv:prop:prov-incons},
    prfND/{pl:ntd:prv:prop:prov-incons},
    prfTab/{pl:tab:prv:prop:prov-incons}}},
له مخې~$\Gamma \Proves !A$۔
\end{proof}

\tagprob{FOL}
\begin{prob}
\olref[fol][com][cth]{cor:completeness} وکاروئ چې
\olref[fol][com][cth]{thm:completeness} ثابت کړئ، او په دې توګه
وښيئ چې د بشپړتيا قضيې دواړه بيانونه هم ارزښته دي۔
\end{prob}
\tagendprob

\tagprob{notFOL}
\begin{prob}
\olref[pl][com][cth]{cor:completeness} وکاروئ چې
\olref[pl][com][cth]{thm:completeness} ثابت کړئ، او په دې توګه
وښيئ چې د بشپړتيا قضيې دواړه بيانونه هم ارزښته دي۔
\end{prob}
\tagendprob

\tagprob{FOL}
\begin{prob}
د دې دپاره چې د~!!a{derivation} نظام بشپړ وي، قاعدې ئې بايد دومره
پياوړې وي چې هر د صدق وړ نۀ سټ ناسازګار ثابت کړي۔ د بشپړتيا د ثبوت
دپاره د~!!{derivation} کومې قاعدې اړينې وې؟ آيا له دې قاعدو څخه
کومه يوه د ثبوت په هېڅ ځاے کښې نۀ ده کارول شوې؟ د دې پوښتنو د ځواب
دپاره داسې يو لړ يا شکل جوړ کړئ چې وښيي د~!!{derivation} کومې
قاعدې په هغو کومو پايلو کښې کارول شوې چې د
\olref[fol][com][cth]{thm:completeness} ثبوت ته رسېږي۔ په دې ثبوتونو
کښې د قاعدو هر ضمني استعمال هم خامخا ياد کړئ۔
\end{prob}
\tagendprob

\tagprob{notFOL}
\begin{prob}
د دې دپاره چې د~!!a{derivation} نظام بشپړ وي، قاعدې ئې بايد دومره
پياوړې وي چې هر د صدق وړ نۀ سټ ناسازګار ثابت کړي۔ د بشپړتيا د ثبوت
دپاره د~!!{derivation} کومې قاعدې اړينې وې؟ آيا له دې قاعدو څخه
کومه يوه د ثبوت په هېڅ ځاے کښې نۀ ده کارول شوې؟ د دې پوښتنو د ځواب
دپاره داسې يو لړ يا شکل جوړ کړئ چې وښيي د~!!{derivation} کومې
قاعدې په هغو کومو پايلو کښې کارول شوې چې د
\olref[pl][com][cth]{thm:completeness} ثبوت ته رسېږي۔ په دې ثبوتونو
کښې د قاعدو هر ضمني استعمال هم خامخا ياد کړئ۔
\end{prob}
\tagendprob

\end{document}
